% ============================================================
%  chapters/ch25-semantic-infrastructure.tex
% ============================================================

\chapter{Semantic Infrastructure}

\chapterepigraph{%
  Roads move matter.\\
  Networks move signals.\\
  Semantic infrastructure moves meaning.
}{Flyxion, \textit{Semantic Infrastructure}}

\sidenote{App.\,M for the\\full category-\\theoretic treatment}

Part~VII begins. The previous parts developed the individual agent
(Part~VI) and the physical universe (Part~V). Part~VII asks: what happens
when agents compose into civilizations? What is the infrastructure through
which meaning persists, accumulates, and evolves across generations?

The answer requires going beyond both the individual agent framework of
Part~VI and the field-theoretic framework of Part~V. The appropriate
language is category theory: the study of composition.

\section{Beyond Files and Folders}

Contemporary digital infrastructure organizes information as files in
folders — a spatial metaphor borrowed from physical paper filing.
This metaphor was adequate for documents but is inadequate for knowledge.

The limitations:

\begin{itemize}
  \item \textbf{No semantics:} a file system knows the name and location
        of a document but nothing about its meaning. Documents about the
        same topic can be anywhere; documents in the same folder can be
        about anything.
  \item \textbf{No composition:} files are concatenated, not composed.
        Merging two research papers produces a combined file, not a
        synthesized understanding.
  \item \textbf{No versioning with meaning:} version control tracks
        character-level changes but not conceptual evolution. A theorem
        that is refactored has changed meaning; a theorem that is
        merely reformatted has not. File-level diffs cannot distinguish
        these.
  \item \textbf{No provenance:} a fact's reliability depends on how it
        was established and by whom. File systems record creation dates
        but not epistemic genealogy.
\end{itemize}

\begin{definition}{Semantic Module}{semantic-module}
  A \emph{semantic module} is a triple $(C, A, \partial C)$ where:
  \begin{itemize}
    \item $C$ is the \emph{content} (propositions, procedures, data),
    \item $A$ is the \emph{admissibility structure} (what inferences
          from $C$ are valid, what modifications preserve meaning),
    \item $\partial C$ is the \emph{interface} (what the module exposes
          to other modules — its public propositions and transformation
          signatures).
  \end{itemize}
\end{definition}

Semantic modules compose through their interfaces. Two modules can be
merged if their interfaces are compatible. The result of merging is a
new module whose content is the synthesis of both, whose admissibility
structure is the intersection of both, and whose interface is the
union of both.

\section{Semantic Modules and Homotopy Merges}

The categorical framework of Appendix~M formalizes composition as
pushouts. But for knowledge, we need a richer notion: \emph{homotopy
merges}, which identify not just syntactically identical structures
but semantically equivalent ones.

\begin{definition}{Homotopy Merge}{homotopy-merge}
  Two semantic modules $(C_1, A_1, \partial C_1)$ and $(C_2, A_2,
  \partial C_2)$ admit a \emph{homotopy merge} if there exists a
  semantic equivalence $\phi : \partial C_1 \to \partial C_2$
  (a homotopy equivalence in the category $\mathbf{Sem}$) between
  their interfaces. The merged module is the pushout:
  \[
    (C_1, A_1, \partial C_1) \sqcup_\phi (C_2, A_2, \partial C_2)
  \]
  with interface $\partial C_1 \cup_\phi \partial C_2$.
\end{definition}

Homotopy merges allow modules developed independently to be combined
even when their interfaces use different terminology or notation,
provided there is a semantic translation between them. This is how
interdisciplinary synthesis works: mathematics developed for physics
is merged with mathematics developed for biology through a homotopy
equivalence (a translation) between the two terminological systems.

\begin{example}
  The module of thermodynamics (entropy, temperature, heat flow) and
  the module of information theory (Shannon entropy, mutual information,
  channel capacity) admit a homotopy merge through the Boltzmann–Shannon
  correspondence: $S = k_B \log W$ translates between the two interface
  languages. The merged module contains both physical and informational
  interpretations of entropy, allowing results from one domain to
  transfer to the other.
\end{example}

\section{Knowledge Evolution}

Knowledge is not static. It evolves through the following operations
on the category $\mathbf{Sem}$:

\begin{itemize}
  \item \textbf{Extension:} adding new propositions to a module's content
        while preserving its interface.
  \item \textbf{Correction:} removing or modifying propositions that have
        been found false, updating the admissibility structure.
  \item \textbf{Composition:} merging two modules (homotopy merge).
  \item \textbf{Abstraction:} forming a quotient module that exposes
        only part of the interface while preserving the admissibility
        structure.
  \item \textbf{Instantiation:} binding abstract modules to specific cases.
\end{itemize}

These operations define a dynamics on $\mathbf{Sem}$ — a flow through
the space of semantic modules. The RSVP accessibility field of
Part~V can be interpreted over this space: high-accessibility regions
contain modules with many potential extensions and compositions;
low-accessibility regions contain modules that are nearly complete
or dead ends.

\section{Semantic Version Control}

Traditional version control (Git, SVN) tracks character-level changes to
files. Semantic version control tracks conceptual-level changes to modules.

\begin{definition}{Semantic Version History}{sem-version}
  A \emph{semantic version history} of module $M$ is a sequence
  \[
    M_0 \xrightarrow{e_1} M_1 \xrightarrow{e_2} M_2 \xrightarrow{e_3}
    \cdots \xrightarrow{e_n} M_n
  \]
  where each $e_i$ is one of the operations (extension, correction,
  composition, abstraction, instantiation) and each $M_i$ is a valid
  semantic module. The \emph{semantic diff} between $M_i$ and $M_j$
  is the minimal sequence of operations connecting them.
\end{definition}

Semantic version control makes knowledge evolution explicit and
reversible. When a correction is made (a proposition removed), the
version history records why and by whom. When modules are composed,
the history records the homotopy equivalence used. The epistemic
genealogy of any proposition is recoverable from the version history.

\section{Obstruction Cohomology}

Not all merges succeed. When two modules cannot be merged — when their
admissibility structures are incompatible — there is an obstruction.
The mathematical theory of obstructions is cohomology (Appendix~I).

\begin{definition}{Semantic Obstruction}{sem-obstruction}
  Two semantic modules $(C_1, A_1, \partial C_1)$ and $(C_2, A_2,
  \partial C_2)$ have a \emph{semantic obstruction} if their
  admissibility structures contradict each other: there exist
  propositions $p_1 \in C_1$ and $p_2 \in C_2$ such that
  $A_1 \models \neg p_2$ and $A_2 \models \neg p_1$.
\end{definition}

Obstructions correspond to genuine contradictions between knowledge
systems — not mere terminological differences (which can be resolved by
homotopy equivalence) but deep incompatibilities. The Copernican
revolution involved an obstruction between the geocentric and
heliocentric modules: the admissibility structure of the geocentric
module (celestial spheres, divine ordering) entailed the falsity
of the heliocentric module's central propositions.

\begin{namedthm}{The Infrastructure Principle}
  Semantic infrastructure is the engineering discipline that builds,
  maintains, and extends the category $\mathbf{Sem}$ — the system
  of semantic modules, their composition rules, their version histories,
  and their obstruction structures. Its goal is to maximize the
  accessibility entropy of the civilizational knowledge base:
  \[
    S_{\mathbf{Sem}} = \log \bigl|\{\text{admissible extensions of }
    \mathbf{Sem}\}\bigr|.
  \]
  A civilization with high $S_{\mathbf{Sem}}$ has a flexible,
  extensible knowledge base; one with low $S_{\mathbf{Sem}}$ has
  a rigid, constrained one.
\end{namedthm}

\begin{exercise}
  Define a metric on $\mathbf{Sem}$ (the category of semantic modules)
  that measures how far two modules are from being mergeable. What
  should such a metric measure? (Consider: number of obstruction pairs,
  complexity of the homotopy equivalence needed, size of the interface
  mismatch.)
\end{exercise}

\begin{exercise}
  The Internet is often described as an information infrastructure.
  In what ways does it fail to be a \emph{semantic} infrastructure?
  Which of the five knowledge evolution operations does it support
  well, and which does it support poorly?
\end{exercise}

\begin{exercise}[title={Constructive}]
  Design a minimal semantic version control system for a single academic
  paper. What constitutes a ``version''? What operations are allowed?
  How would ``correction'' and ``composition'' be represented in
  your system? What would the semantic diff look like?
\end{exercise}
